\documentclass[11pt,a4paper]{article}
\usepackage[margin=2.2cm]{geometry}
\usepackage{booktabs,array,longtable}
\usepackage[table]{xcolor}
\usepackage{titlesec}
\usepackage{enumitem}
\usepackage[colorlinks,allcolors=blue!60!black]{hyperref}
\usepackage{parskip}
\usepackage{amssymb}
\definecolor{hd}{RGB}{20,50,90}
\titleformat{\section}{\Large\bfseries\color{hd}}{\thesection}{0.6em}{}
\titleformat{\subsection}{\large\bfseries\color{hd}}{\thesubsection}{0.6em}{}
\newcommand{\yes}{\textcolor{green!50!black}{\textbf{have}}}
\newcommand{\no}{\textcolor{red!70!black}{\textbf{missing}}}
\newcommand{\parti}{\textcolor{orange!80!black}{\textbf{partial}}}
\setlist{itemsep=2pt,topsep=3pt,parsep=0pt}

\title{\vspace{-1.5cm}\textbf{COMPASS at CVPR}\\[2mm]
\large What the results actually say, what the dataset must become,\\and what to collect between now and 16 November 2026}
\author{Planning document}
\date{8 September 2026}

\begin{document}
\maketitle
\vspace{-6mm}
\hrule
\vspace{4mm}

\section{Correction: we are stronger than the README says}

The repository README reports COMPASS-WNet at $9.97$~dB against RadioUNet at $10.00$~dB
over 120 samples and calls it a tie. That table is stale. The current
\texttt{results/dl\_baselines.json}, over 320 samples with bootstrap intervals, says
something materially different.

\begin{center}
\small
\begin{tabular}{lrrl}
\toprule
Method & Free-unobs RMSE (dB) & SSIM & Uncertainty corr. \\
\midrule
\rowcolor{green!8}\textbf{COMPASS-WNet + occlusion} & \textbf{8.65} \small[8.42,\,8.86] & \textbf{0.841} & 0.553 \\
RadioUNet & 9.93 \small[9.68,\,10.18] & 0.816 & none \\
COMPASS-WNet (no occlusion) & 10.11 \small[9.82,\,10.42] & 0.823 & 0.530 \\
RadioMamba & 10.92 & 0.795 & 0.486 \\
RadioTransformer & 10.97 & 0.791 & 0.454 \\
URAM & 11.34 & 0.800 & 0.436 \\
RadioDiff & 11.40 & 0.802 & none \\
RadioGAN & 11.42 & 0.788 & none \\
PMNet & 14.12 & 0.714 & none \\
SparseUNet (no geometry) & 21.05 & 0.656 & $-0.022$ \\
RBF / Ordinary Kriging / IDW / GP & 25.6 to 28.1 & 0.67 to 0.72 & 0.34 to 0.38 \\
RMDM & 35.98 & 0.715 & none \\
\bottomrule
\end{tabular}
\end{center}

Three things follow, and they change my earlier advice.

\textbf{You beat the field outright.} $8.65$ against RadioUNet's $9.93$ is a $12.9\%$
reduction, and the $95\%$ intervals $[8.42, 8.86]$ and $[9.68, 10.18]$ do not overlap.
This is not a tie. Against the best classical method it is a factor of three. My earlier
concern that CVPR rewards clear wins and you only had a tie is void.

\textbf{The win is attributable to one idea.} COMPASS-WNet scores $10.11$ without the
occlusion field and $8.65$ with it. The single geometric channel is worth $1.46$~dB,
which is larger than the entire gap to the next best method. That is the cleanest
possible ablation story: one physically motivated input, one interpretable mechanism,
one decisive number.

\textbf{You own uncertainty.} Every competitive baseline reports none. Among those that
do, yours correlates best with actual error at $0.553$.

\section{What is still missing for CVPR}

The synthetic story is now strong enough. Two gaps remain, and only one of them is
serious.

\textbf{Architectural novelty is modest, and that is survivable.} A reviewer will see a
cascaded U-Net with extra input channels. The defence is that the contribution is not
the backbone but the diagnostic and the channel: unrelated architectures plateau at the
\emph{same} non-line-of-sight error, which identifies a missing input rather than
insufficient capacity, and the occlusion field supplies exactly the line integral a
convolution cannot form. Framed that way it reads as insight rather than engineering.
Many accepted CVPR papers are one good idea attached to a standard backbone.

\textbf{The real-data result is the actual problem.} On UniCellular the learned
reconstructor reaches $6.30$~dB after fine-tuning while classical RBF reaches
$5.51$~dB, and \texttt{sim2real\_finetune.json} records
\texttt{finetuned\_beats\_classical: false}. A residual model on top of classical does
edge ahead at $5.14$ against $5.49$, but the intervals overlap. As it stands, the paper
must either omit real data, which wastes its best differentiator, or report that
interpolation wins, which a reviewer reads as a failed method.

\subsection{Why the real-data result is a data problem, not a method problem}

This is the key diagnosis. Real evaluation currently uses held-out reference points from
54 to 67 usable cells, roughly 2{,}000 evaluation points. There is no dense ground truth
anywhere in the corpus, so you cannot measure reconstruction error over a field. You can
only measure error at sparse survey points that were themselves collected the same way
as the training points. Under that protocol an interpolator is close to optimal by
construction, because the metric only ever asks it to interpolate between nearby samples
of the same process. The learned model's advantages, which are extrapolation into
unmeasured regions and geometry-aware behaviour behind walls, are precisely the things
the protocol cannot see.

Fixing the protocol, not the model, is what turns this section from a liability into the
paper's headline.

\section{The proposed headline and storyline}

\subsection{Headline}

\begin{quote}
\textbf{Simulation says radio-map reconstruction is solved. On real measurements with
dense ground truth, every method including ours degrades sharply, and the ranking
changes. We provide the benchmark that shows it and the geometric prior that narrows the
gap.}
\end{quote}

This is a shape CVPR recognises and rewards, the ``in the wild'' paper. It converts your
current weakness into the contribution. You are no longer a method that loses to
interpolation on real data. You are the group that built the first real benchmark on
which the question can be asked honestly, and you happen to hold the best method on both
halves.

\subsection{Narrative arc}

\begin{enumerate}
\item Radio-map reconstruction is a sparse-to-dense inverse problem whose field is
governed by geometry that is directly observable. Vision should care.
\item On the standard synthetic benchmark the problem looks close to solved, and we push
it further with a ray-occlusion field, reaching $8.65$~dB against $9.93$ for the best
prior method.
\item We then show the benchmark is measuring the wrong thing. We introduce a real
multi-site, multi-floor indoor benchmark with dense ground truth, trajectories, inertial
signals and surveyed transmitters.
\item On it, learned methods that dominate simulation fall behind classical
interpolation, and the ordering of published methods does not survive transfer. We
characterise the gap along geometry, device heterogeneity and acquisition order.
\item The geometric prior transfers where the rest does not, and calibrated uncertainty
turns into a usable acquisition rule.
\end{enumerate}

\subsection{Candidate titles}

\begin{itemize}
\item \textbf{Walls Cast Soft Shadows: Geometry-Conditioned Radio Field Reconstruction
in Simulation and Reality}
\item \textbf{The Sim-to-Real Gap in Radio Map Reconstruction}
\item \textbf{Penumbra: Occlusion-Aware Radio Field Reconstruction with a Real-World
Benchmark}
\end{itemize}

On naming, \textsc{compass} is heavily overloaded in machine learning.
\textbf{Penumbra} encodes the finding that walls attenuate rather than block, it is
native vision and graphics vocabulary, and it is distinctive. The draft defines
\texttt{\textbackslash methodname} once so the change costs one line.

\section{The dataset contribution}

\subsection{Hard constraint}

CVPR reviewer guidelines are explicit. A dataset may be \emph{used} while private, but it
may only be \emph{claimed as a contribution} if there is a reasonable expectation of
public release on publication. If the data cannot be released, drop the dataset claim and
keep the method claim. This decision gates everything below and should be settled first,
because it involves the building owners at Ezdan and the CMUQ estates office, not us.

\subsection{What already exists}

\begin{center}
\small
\begin{tabular}{lrrrrr}
\toprule
Site & Floors & Phones & Transmitters & Stationary rows & Mobile rows \\
\midrule
CMUQ & 3 & 7 & 57 & 2{,}174{,}421 & 729{,}207 \\
EC Parking & 3 & 6 & 41 & 628{,}732 & 328{,}787 \\
Ezdan Tower 1 & 30 & 6 & 64 & 531{,}863 & 210{,}248 \\
Ezdan Tower 4 & 25 & 6 & 42 & 395{,}103 & 66{,}369 \\
\midrule
\textbf{Total} & \textbf{61} & 6 to 7 & \textbf{204} & \textbf{3{,}730{,}119} & \textbf{1{,}334{,}611} \\
\bottomrule
\end{tabular}
\end{center}

This is already a serious corpus: four buildings, sixty-one floors including two
high-rises, two hundred transmitters, five million rows, seven devices. Scale is not the
problem. Per-sample schema is
\texttt{phoneName}, \texttt{rpNumber}, \texttt{scanNumber}, \texttt{timeStamp},
\texttt{transmitter\_id}, \texttt{transmitter\_rss}, \texttt{coordinates},
\texttt{floor}, plus floor plans.

\subsection{What is missing, in priority order}

\begin{center}
\small
\renewcommand{\arraystretch}{1.25}
\begin{tabular}{>{\raggedright}p{4.3cm}cp{8.6cm}}
\toprule
Quantity & Status & Why it matters and what it unlocks \\
\midrule
Dense ground-truth field & \no &
\textbf{The single blocking gap.} Without a densely sampled reference on at least one
floor you cannot measure reconstruction error, only interpolation error at survey
points. This is why classical wins today. Fixing it is what makes the real benchmark
real. \\
Surveyed transmitter positions & \parti &
Only $18\%$ of cells currently yield a reliable point source. Ground-truth antenna and
booster locations for a subset let you separate ``the model cannot localise'' from ``the
model cannot propagate'', and make the occlusion field computable on real data. \\
Continuous trajectory ground truth & \parti &
\texttt{rpNumber} implies discrete survey points. Mobile rows exist with lag-1
autocorrelation $0.77$ to $0.97$, but without accurate continuous position the order and
active-sensing claims rest on weak positioning. \\
Inertial signals (IMU) & \no &
Accelerometer, gyroscope, magnetometer and barometer. Turns this into a multimodal
sensing dataset, supports dead reckoning between fixes, makes the order argument
mechanistic rather than statistical, and is what a vision audience recognises as a
modern dataset. \\
Device calibration references & \parti &
Offsets are estimable at $0.78$ to $0.89$ correlation, but a controlled
same-place-same-time capture across all seven phones gives a ground-truth offset per
device rather than an inferred one. \\
Radio metadata & \parti &
Band, frequency, PCI, cell type, indoor or outdoor. Needed to explain why boosters share
identifiers and to justify the transmitter-agnostic design. \\
Temporal repeats & \no &
The same route walked on different days, times and crowd levels. Quantifies drift, and
gives a generalisation axis reviewers will ask about. \\
Orientation and body effects & \no &
Phone held, pocketed, bag. Body attenuation is several dB and is an obvious confound. \\
\bottomrule
\end{tabular}
\end{center}

\section{Collection plan}

Resources assumed: six to seven phones, a team of people, access to CMUQ and the existing
sites, roughly ten weeks to 16 November 2026.

\subsection{Priority 1: the dense ground-truth floor (non-negotiable)}

Pick \textbf{one} floor. CMUQ floor 2 is the natural choice because it already has the
most transmitters, the widest device spread at $13.5$~dB and six shared phones.

\begin{itemize}
\item Survey on a regular grid, target $1$ to $2$~m spacing across all accessible free
space. For a floor of roughly $80\times60$~m at $2$~m spacing this is on the order of
1{,}200 points.
\item At each point, dwell for a fixed window, for example 30 seconds, and record every
visible cell on every phone simultaneously.
\item Mark the point physically or use a total station, laser rangefinder or a
tape-and-chalk grid. Positional error should be well under the grid pitch.
\item This yields, for the first time, a field you can hold out in blocks rather than at
scattered points. \textbf{Evaluate by masking contiguous regions}, which is what
finally lets extrapolation and geometry show their value.
\end{itemize}

Estimated effort: two people, three to four days per floor at 2~m spacing, plus a day of
setup. This is the highest-value use of the available time by a wide margin.

\subsection{Priority 2: trajectories with usable position}

\begin{itemize}
\item Walk repeatable routes with start, end and periodic waypoints marked on the floor
plan, and press a hardware or on-screen marker at each waypoint so position can be
interpolated between known fixes rather than guessed.
\item Log continuously at the highest scan rate the handset allows, and record the
scan-to-timestamp mapping.
\item Cover at least: a corridor loop, a route crossing several walls, a route into a
deep non-line-of-sight region, and an open-space route as a control.
\item Repeat each route with at least three phones carried together, and at three
distinct times of day.
\end{itemize}

\subsection{Priority 3: inertial and context streams}

Record alongside every mobile session, at the stated minimum rates:

\begin{itemize}
\item Accelerometer, gyroscope, magnetometer at $\ge 50$~Hz.
\item Barometer at $\ge 1$~Hz for floor transitions.
\item Step counter and device orientation quaternion if exposed.
\item Screen state, and whether the phone was in hand, pocket or bag, logged per session.
\item A monotonic clock alongside wall-clock time, so drift between devices is
recoverable.
\end{itemize}

\subsection{Priority 4: transmitter survey}

\begin{itemize}
\item Physically locate and record the position of every accessible antenna, femtocell
and booster on the dense floor, with floor and mounting height.
\item Record the mapping from physical unit to the cell identifiers it broadcasts, which
is what resolves the shared-identifier problem.
\item Even twenty surveyed units on one floor changes the transmitter-agnostic claim from
an assumption into a measured statement.
\end{itemize}

\subsection{Priority 5: device calibration capture}

\begin{itemize}
\item All seven phones, same location, same orientation, same two-minute window,
repeated at five positions spanning strong to weak signal.
\item Repeat at the start and end of the campaign to capture drift.
\item Gives a ground-truth per-device offset curve as a function of signal level, rather
than a single scalar.
\end{itemize}

\subsection{Clock synchronisation}

Cross-device order and device comparison both depend on it. Synchronise all phones to a
single NTP source at the start of each session, log both wall-clock and monotonic time,
and begin every session with a shared physical event, for example all phones held
together for ten seconds, so residual offset is recoverable in post-processing.

\subsection{Privacy and release}

Because release gates the dataset claim, handle it at collection time rather than after.
\begin{itemize}
\item Collect with team members as the only participants, so there are no third-party
subjects and no consent chain to reconstruct later.
\item Record no personally identifying information. Replace \texttt{phoneName} with a
stable pseudonym at capture.
\item Obtain written permission from CMUQ and Ezdan for floor plans and cell identifiers
to be published. Floor plans are the item most likely to be refused, so ask early and be
ready to release a simplified occupancy raster instead of an architectural drawing.
\item Decide whether transmitter positions can be published. If not, release the derived
occlusion field rather than the coordinates.
\end{itemize}

\section{What the real-data headline becomes}

With the dense floor in hand, the section rewrites itself along these lines.

\begin{itemize}
\item \textbf{The ranking does not transfer.} Methods are re-evaluated on a real dense
field with block-held-out regions. Report the reordering explicitly. This is the
scientific payload and is publishable regardless of who wins.
\item \textbf{The gap is quantified.} Same architectures, same protocol, simulation
against reality, with the degradation attributed across geometry, device heterogeneity
and sampling structure.
\item \textbf{Geometry transfers.} The occlusion field is computed from a floor plan and
a surveyed transmitter, so it is available on real data. If the $1.46$~dB it is worth in
simulation survives even partially, that is the paper's strongest single sentence.
\item \textbf{Uncertainty becomes actionable.} With a dense field you can finally show
that guided acquisition reaches a target error with fewer measurements than
space-filling, on real data, which is the practical claim that makes the work matter.
\end{itemize}

\section{Ten weeks to the deadline}

\begin{center}
\small
\begin{tabular}{llp{9.2cm}}
\toprule
Weeks & Focus & Detail \\
\midrule
1 & Decide and unblock & Settle the release question with CMUQ and Ezdan. Choose the
dense floor. Build and test the logging app with IMU and marker support. Pilot on one
corridor and verify every stream lands. \\
2 to 4 & Collect & Dense grid survey. Trajectory routes with repeats. Transmitter survey.
Device calibration captures. \\
4 to 6 & Process and release-prep & Position interpolation, clock alignment, occupancy
rasters, per-device offsets, anonymisation, dataset card and licence. \\
5 to 8 & Experiments & Re-run every baseline under the block-held-out protocol. Sim-to-real
transfer matrix. Occlusion ablation on real data. Guided acquisition on the dense field. \\
8 to 10 & Write and buffer & Full draft by week 9. Keep a week of slack, since collection
always overruns. \\
\bottomrule
\end{tabular}
\end{center}

\section{Risks, and the honest fallback}

\begin{itemize}
\item \textbf{Release refused.} Most likely failure. Fallback is to drop the dataset
claim, keep the real-data evaluation as supporting analysis, and submit on the method
plus the sim-to-real characterisation. Still a viable paper.
\item \textbf{Dense survey does not finish.} Reduce scope to a partial floor rather than
coarsening the grid. A small dense region is scientifically useful, a large sparse one is
not, because sparse is exactly the regime that already fails to discriminate.
\item \textbf{Geometry does not transfer.} Genuinely possible, and it must be reported if
so. It remains publishable as a negative transfer result on a new benchmark, which is
why the benchmark should be the primary claim and the method the secondary one.
\item \textbf{Ten weeks is tight.} If the collection slips past week 6, target a
sensing venue where the contribution is legible without the vision framing, and use the
CVPR draft as the base.
\end{itemize}

\section{Decisions needed from you}

\begin{enumerate}
\item Can the data be released publicly? Everything else follows from this.
\item Which floor becomes the dense reference, and who can commit the days to survey it?
\item Do we rename to Penumbra, or keep COMPASS for continuity with the repository and
prior decks?
\item Is the dataset the primary contribution with the method second, or the reverse? I
recommend the former, because it is the more defensible claim and it makes the real-data
result an asset rather than a liability.
\end{enumerate}

\end{document}
